% Append to Section V-G, "Projection Placement and Computation", after
% the existing timing discussion. Do not replace the older timing numbers.

We also time the six classical methods from
Section~\ref{sec:revision-structured-comparisons} in a separate serial
experiment. This uses the vectorized comparator implementation,
$N=32$, $K=3$, $P=30$, $L=7$, SNR $=5$~dB and RSR $=12$~dB.
Table~\ref{tab:revision-serial} reports medians over 20 fresh worlds with
single-thread CPU execution and double precision. The iterative methods
use 100 final updates, 25 updates per candidate and four structural
sweeps. Timings include initialization, parameter selection where needed,
and the final fit, but exclude channel and observation generation.

\begin{table}[!t]
\caption{Serial runtime of the classical comparator implementations.}
\label{tab:revision-serial}
\centering
\footnotesize
\begin{tabular}{lr}
\toprule
Method & Median total time (ms)\\
\midrule
EM-GS & 9.943\\
HS-GS & 348.536\\
Rician GD & 10.015\\
Rician GD + Hankel & 346.500\\
EM-GS + smooth filtering & 200.799\\
Xu 1D linearized LS & 0.647\\
\bottomrule
\end{tabular}
\end{table}

The methods are timed in a fixed order for each world: HS-GS,
Rician--Hankel, smooth filtering, EM-GS, Rician GD and linearized LS.
No warm-up worlds are discarded. The median selection share is 76.0\%
for HS-GS and 56.4\% for smooth filtering; the latter searches seven
threshold factors rather than 16 ranks. These are measured costs for one
implementation, machine and operating point.
They are separate from the earlier production-code placement timings at
$P=20$ and RSR $=10$~dB, so a difference between the two timing studies
must not be interpreted as an algorithmic speedup.
